% =====================================================================
%  Anatomía de un LLM — Fase III — Módulo 13
%  Atención
%  Compilar: latexmk -r ../latexmkrc M13-atencion.tex
% =====================================================================
\documentclass[11pt,a4paper,twoside,openany]{book}
\usepackage{../estilo-apuntes}

\newcommand{\ApuntesCurso}{Anatomía de un LLM}
\newcommand{\ApuntesModulo}{Módulo 13}

\hypersetup{pdftitle={Módulo 13. Atención (resumen)},pdfauthor={Anatomía de un LLM}}

% Notación propia del módulo
\newcommand{\T}{^{\top}}
\newcommand{\uno}{\mathbf{1}}
\newcommand{\diag}{\operatorname{diag}}
\newcommand{\Var}{\operatorname{Var}}
\newcommand{\Cov}{\operatorname{Cov}}
\newcommand{\rg}{\operatorname{rg}}
\newcommand{\Att}{\operatorname{Attn}}
\newcommand{\MHA}{\operatorname{MHA}}
\newcommand{\Concat}{\operatorname{Concat}}
\newcommand{\lse}{\operatorname{lse}}
\newcommand{\conv}{\operatorname{conv}}
\newcommand{\FLOP}{\mathrm{FLOP}}
\newcommand{\WQ}{W_{Q}}
\newcommand{\WK}{W_{K}}
\newcommand{\WV}{W_{V}}
\newcommand{\WO}{W_{O}}
\newcommand{\tk}[1]{\text{«#1»}}

\makeatletter
\newlabel{fig:qkv}{{1.1}{2}{Autoatención para la posición $i$. Tres proyecciones aprendidas asignan a cada token los papeles de consulta, clave y valor; los parecidos de $q_i$ con todas las claves se normalizan con una softmax y ponderan los valores}{figure.1.1}{}}
\newlabel{fig:qkv@cref}{{[figure][1][1]1.1}{[1][1][]2}}
\newlabel{eje:1-kv}{{1.1}{4}{}{ejercicio.1.1}{}}
\newlabel{eje:1-kv@cref}{{[ejercicio][1][1]1.1}{[1][4][]4}}
\newlabel{eje:1-escalar}{{1.2}{4}{}{ejercicio.1.2}{}}
\newlabel{eje:1-escalar@cref}{{[ejercicio][2][1]1.2}{[1][4][]4}}
\newlabel{eje:1-aditiva}{{1.3}{4}{}{ejercicio.1.3}{}}
\newlabel{eje:1-aditiva@cref}{{[ejercicio][3][1]1.3}{[1][4][]4}}
\newlabel{eje:1-copiar}{{1.4}{4}{}{ejercicio.1.4}{}}
\newlabel{eje:1-copiar@cref}{{[ejercicio][4][1]1.4}{[1][4][]4}}
\newlabel{eje:1-limites}{{1.5}{4}{}{ejercicio.1.5}{}}
\newlabel{eje:1-limites@cref}{{[ejercicio][5][1]1.5}{[1][4][]4}}
\newlabel{ej:tres}{{2.2}{5}{Tres tokens}{ejemplo.2}{}}
\newlabel{ej:tres@cref}{{[ejemplo][2][]2.2}{[1][5][]5}}
\newlabel{fig:mapa}{{2.1}{6}{Mapa de atención del \cref {ej:tres}; la intensidad del gris es proporcional a $\alpha _{ij}$. La fila $i$ indica a qué posiciones atiende la posición $i$}{figure.2.1}{}}
\newlabel{fig:mapa@cref}{{[figure][1][2]2.1}{[1][6][]6}}
\newlabel{eje:2-formas}{{2.1}{8}{}{ejercicio.2.1}{}}
\newlabel{eje:2-formas@cref}{{[ejercicio][1][2]2.1}{[1][8][]8}}
\newlabel{eje:2-var}{{2.2}{8}{}{ejercicio.2.2}{}}
\newlabel{eje:2-var@cref}{{[ejercicio][2][2]2.2}{[1][8][]8}}
\newlabel{eje:2-temp}{{2.3}{8}{}{ejercicio.2.3}{}}
\newlabel{eje:2-temp@cref}{{[ejercicio][3][2]2.3}{[1][8][]8}}
\newlabel{eje:2-jac}{{2.4}{8}{}{ejercicio.2.4}{}}
\newlabel{eje:2-jac@cref}{{[ejercicio][4][2]2.4}{[1][8][]8}}
\newlabel{eje:2-sigma}{{2.5}{8}{}{ejercicio.2.5}{}}
\newlabel{eje:2-sigma@cref}{{[ejercicio][5][2]2.5}{[1][8][]8}}
\newlabel{eje:2-cruzada}{{2.6}{8}{}{ejercicio.2.6}{}}
\newlabel{eje:2-cruzada@cref}{{[ejercicio][6][2]2.6}{[1][8][]8}}
\newlabel{fig:causal}{{3.1}{10}{Máscara causal para $n=5$. La fila $i$ solo permite las posiciones $j\le i$; las celdas sombreadas reciben peso exactamente nulo}{figure.3.1}{}}
\newlabel{fig:causal@cref}{{[figure][1][3]3.1}{[1][9][]10}}
\newlabel{ej:mascara}{{3.5}{10}{}{ejemplo.5}{}}
\newlabel{ej:mascara@cref}{{[ejemplo][5][]3.5}{[1][10][]10}}
\newlabel{eje:3-despues}{{3.1}{11}{}{ejercicio.3.1}{}}
\newlabel{eje:3-despues@cref}{{[ejercicio][1][3]3.1}{[1][11][]11}}
\newlabel{eje:3-jm}{{3.2}{11}{}{ejercicio.3.2}{}}
\newlabel{eje:3-jm@cref}{{[ejercicio][2][3]3.2}{[1][11][]11}}
\newlabel{eje:3-calculo}{{3.3}{11}{}{ejercicio.3.3}{}}
\newlabel{eje:3-calculo@cref}{{[ejercicio][3][3]3.3}{[1][11][]11}}
\newlabel{eje:3-finito}{{3.4}{12}{}{ejercicio.3.4}{}}
\newlabel{eje:3-finito@cref}{{[ejercicio][4][3]3.4}{[1][12][]12}}
\newlabel{eje:3-renorm}{{3.5}{12}{}{ejercicio.3.5}{}}
\newlabel{eje:3-renorm@cref}{{[ejercicio][5][3]3.5}{[1][12][]12}}
\newlabel{fig:mha}{{4.1}{13}{Atención multi-cabeza: $h$ atenciones independientes en subespacios de dimensión $d/h$, concatenadas y mezcladas por $\WO $}{figure.4.1}{}}
\newlabel{fig:mha@cref}{{[figure][1][4]4.1}{[1][13][]13}}
\newlabel{eje:4-vf}{{4.1}{16}{}{ejercicio.4.1}{}}
\newlabel{eje:4-vf@cref}{{[ejercicio][1][4]4.1}{[1][16][]16}}
\newlabel{eje:4-param}{{4.2}{16}{}{ejercicio.4.2}{}}
\newlabel{eje:4-param@cref}{{[ejercicio][2][4]4.2}{[1][16][]16}}
\newlabel{eje:4-flops}{{4.3}{16}{}{ejercicio.4.3}{}}
\newlabel{eje:4-flops@cref}{{[ejercicio][3][4]4.3}{[1][16][]16}}
\newlabel{eje:4-mem}{{4.4}{16}{}{ejercicio.4.4}{}}
\newlabel{eje:4-mem@cref}{{[ejercicio][4][4]4.4}{[1][16][]16}}
\newlabel{eje:4-jm}{{4.5}{16}{}{ejercicio.4.5}{}}
\newlabel{eje:4-jm@cref}{{[ejercicio][5][4]4.5}{[1][16][]16}}
\newlabel{eje:4-cs336}{{4.6}{16}{}{ejercicio.4.6}{}}
\newlabel{eje:4-cs336@cref}{{[ejercicio][6][4]4.6}{[1][16][]16}}
\newlabel{eje:4-rango}{{4.7}{17}{}{ejercicio.4.7}{}}
\newlabel{eje:4-rango@cref}{{[ejercicio][7][4]4.7}{[1][17][]17}}
\newlabel{fig:induccion}{{5.1}{19}{Cabeza de inducción. La capa~1 escribe en la posición~2 que su token anterior es «el»; en una posición posterior $i$ con el token «el», la consulta de la capa~2 encaja con esa clave y el circuito OV copia «gato», que se convierte en la predicción para la posición $i+1$}{figure.5.1}{}}
\newlabel{fig:induccion@cref}{{[figure][1][5]5.1}{[1][19][]19}}
\newlabel{eje:5-capa}{{5.1}{21}{}{ejercicio.5.1}{}}
\newlabel{eje:5-capa@cref}{{[ejercicio][1][5]5.1}{[1][21][]21}}
\newlabel{eje:5-previo}{{5.2}{21}{}{ejercicio.5.2}{}}
\newlabel{eje:5-previo@cref}{{[ejercicio][2][5]5.2}{[1][21][]21}}
\newlabel{eje:5-hop}{{5.3}{21}{}{ejercicio.5.3}{}}
\newlabel{eje:5-hop@cref}{{[ejercicio][3][5]5.3}{[1][21][]21}}
\newlabel{eje:5-induccion}{{5.4}{21}{}{ejercicio.5.4}{}}
\newlabel{eje:5-induccion@cref}{{[ejercicio][4][5]5.4}{[1][21][]21}}
\newlabel{eje:5-jac}{{5.5}{21}{}{ejercicio.5.5}{}}
\newlabel{eje:5-jac@cref}{{[ejercicio][5][5]5.5}{[1][21][]21}}
\newlabel{eje:5-recup}{{5.6}{21}{}{ejercicio.5.6}{}}
\newlabel{eje:5-recup@cref}{{[ejercicio][6][5]5.6}{[1][21][]21}}
\newlabel{ap:soluciones}{{A}{22}{Soluciones de los ejercicios}{appendix.A}{}}
\newlabel{ap:soluciones@cref}{{[appendix][1][2147483647]A}{[1][22][]22}}
\newlabel{sol:eje:1-kv}{{A.1}{22}{Tema 1}{section*.3}{}}
\newlabel{sol:eje:1-kv@cref}{{[subappendix][1][2147483647,1]A.1}{[1][22][]22}}
\newlabel{sol:eje:1-escalar}{{A.1}{22}{Tema 1}{section*.4}{}}
\newlabel{sol:eje:1-escalar@cref}{{[subappendix][1][2147483647,1]A.1}{[1][22][]22}}
\newlabel{sol:eje:1-aditiva}{{A.1}{23}{Tema 1}{section*.5}{}}
\newlabel{sol:eje:1-aditiva@cref}{{[subappendix][1][2147483647,1]A.1}{[1][23][]23}}
\newlabel{sol:eje:1-copiar}{{A.1}{23}{Tema 1}{section*.6}{}}
\newlabel{sol:eje:1-copiar@cref}{{[subappendix][1][2147483647,1]A.1}{[1][23][]23}}
\newlabel{sol:eje:1-limites}{{A.1}{24}{Tema 1}{section*.7}{}}
\newlabel{sol:eje:1-limites@cref}{{[subappendix][1][2147483647,1]A.1}{[1][24][]24}}
\newlabel{sol:eje:2-formas}{{A.2}{24}{Tema 2}{section*.8}{}}
\newlabel{sol:eje:2-formas@cref}{{[subappendix][2][2147483647,1]A.2}{[1][24][]24}}
\newlabel{sol:eje:2-var}{{A.2}{25}{Tema 2}{section*.9}{}}
\newlabel{sol:eje:2-var@cref}{{[subappendix][2][2147483647,1]A.2}{[1][25][]25}}
\newlabel{sol:eje:2-temp}{{A.2}{25}{Tema 2}{section*.10}{}}
\newlabel{sol:eje:2-temp@cref}{{[subappendix][2][2147483647,1]A.2}{[1][25][]25}}
\newlabel{sol:eje:2-jac}{{A.2}{25}{Tema 2}{section*.11}{}}
\newlabel{sol:eje:2-jac@cref}{{[subappendix][2][2147483647,1]A.2}{[1][25][]25}}
\newlabel{sol:eje:2-sigma}{{A.2}{26}{Tema 2}{section*.12}{}}
\newlabel{sol:eje:2-sigma@cref}{{[subappendix][2][2147483647,1]A.2}{[1][26][]26}}
\newlabel{sol:eje:2-cruzada}{{A.2}{26}{Tema 2}{section*.13}{}}
\newlabel{sol:eje:2-cruzada@cref}{{[subappendix][2][2147483647,1]A.2}{[1][26][]26}}
\newlabel{sol:eje:3-despues}{{A.3}{27}{Tema 3}{section*.14}{}}
\newlabel{sol:eje:3-despues@cref}{{[subappendix][3][2147483647,1]A.3}{[1][27][]27}}
\newlabel{sol:eje:3-jm}{{A.3}{27}{Tema 3}{section*.15}{}}
\newlabel{sol:eje:3-jm@cref}{{[subappendix][3][2147483647,1]A.3}{[1][27][]27}}
\newlabel{sol:eje:3-calculo}{{A.3}{27}{Tema 3}{section*.16}{}}
\newlabel{sol:eje:3-calculo@cref}{{[subappendix][3][2147483647,1]A.3}{[1][27][]27}}
\newlabel{sol:eje:3-finito}{{A.3}{28}{Tema 3}{section*.17}{}}
\newlabel{sol:eje:3-finito@cref}{{[subappendix][3][2147483647,1]A.3}{[1][28][]28}}
\newlabel{sol:eje:3-renorm}{{A.3}{28}{Tema 3}{section*.18}{}}
\newlabel{sol:eje:3-renorm@cref}{{[subappendix][3][2147483647,1]A.3}{[1][28][]28}}
\newlabel{sol:eje:4-vf}{{A.4}{29}{Tema 4}{section*.19}{}}
\newlabel{sol:eje:4-vf@cref}{{[subappendix][4][2147483647,1]A.4}{[1][29][]29}}
\newlabel{sol:eje:4-param}{{A.4}{29}{Tema 4}{section*.20}{}}
\newlabel{sol:eje:4-param@cref}{{[subappendix][4][2147483647,1]A.4}{[1][29][]29}}
\newlabel{sol:eje:4-flops}{{A.4}{29}{Tema 4}{section*.21}{}}
\newlabel{sol:eje:4-flops@cref}{{[subappendix][4][2147483647,1]A.4}{[1][29][]29}}
\newlabel{sol:eje:4-mem}{{A.4}{30}{Tema 4}{section*.22}{}}
\newlabel{sol:eje:4-mem@cref}{{[subappendix][4][2147483647,1]A.4}{[1][30][]30}}
\newlabel{sol:eje:4-jm}{{A.4}{30}{Tema 4}{section*.23}{}}
\newlabel{sol:eje:4-jm@cref}{{[subappendix][4][2147483647,1]A.4}{[1][30][]30}}
\newlabel{sol:eje:4-cs336}{{A.4}{30}{Tema 4}{section*.24}{}}
\newlabel{sol:eje:4-cs336@cref}{{[subappendix][4][2147483647,1]A.4}{[1][30][]30}}
\newlabel{sol:eje:4-rango}{{A.4}{31}{Tema 4}{section*.25}{}}
\newlabel{sol:eje:4-rango@cref}{{[subappendix][4][2147483647,1]A.4}{[1][31][]31}}
\newlabel{sol:eje:5-capa}{{A.5}{32}{Tema 5}{section*.26}{}}
\newlabel{sol:eje:5-capa@cref}{{[subappendix][5][2147483647,1]A.5}{[1][32][]32}}
\newlabel{sol:eje:5-previo}{{A.5}{32}{Tema 5}{section*.27}{}}
\newlabel{sol:eje:5-previo@cref}{{[subappendix][5][2147483647,1]A.5}{[1][32][]32}}
\newlabel{sol:eje:5-hop}{{A.5}{33}{Tema 5}{section*.28}{}}
\newlabel{sol:eje:5-hop@cref}{{[subappendix][5][2147483647,1]A.5}{[1][32][]33}}
\newlabel{sol:eje:5-induccion}{{A.5}{33}{Tema 5}{section*.29}{}}
\newlabel{sol:eje:5-induccion@cref}{{[subappendix][5][2147483647,1]A.5}{[1][33][]33}}
\newlabel{sol:eje:5-jac}{{A.5}{33}{Tema 5}{section*.30}{}}
\newlabel{sol:eje:5-jac@cref}{{[subappendix][5][2147483647,1]A.5}{[1][33][]33}}
\newlabel{sol:eje:5-recup}{{A.5}{34}{Tema 5}{section*.31}{}}
\newlabel{sol:eje:5-recup@cref}{{[subappendix][5][2147483647,1]A.5}{[1][34][]34}}
\makeatother

\begin{document}
\frontmatter
\begin{titlepage}
\centering
\vspace*{2.5cm}
{\large\scshape Anatomía de un LLM}\par
\vspace{0.4cm}
{\normalsize Fundamentos matemáticos de los modelos de lenguaje}\par
\vspace{2.2cm}
\rule{\textwidth}{0.8pt}\par\vspace{0.5cm}
{\small\scshape Fase III \textperiodcentered\ Secuencias y Transformer}\par\vspace{0.3cm}
{\Huge\bfseries Módulo 13\par}\vspace{0.3cm}
{\LARGE Atención\par}
\vspace{0.5cm}\rule{\textwidth}{0.4pt}\par
\vspace{1.6cm}
{\large Resumen del módulo}\par
\vfill
{\small Edición de octubre de 2026}\par
\end{titlepage}

\chapter*{Sobre este resumen}
\addcontentsline{toc}{chapter}{Sobre este resumen}
\markboth{Sobre este resumen}{}

Este documento reúne lo indispensable de los apuntes del Módulo 13: el texto expositivo, todas las definiciones y todos los resultados (teoremas, proposiciones, lemas y corolarios), las tablas y las observaciones de contenido. Se han suprimido las demostraciones, los ejemplos, las lecturas, las figuras, los ejercicios y sus soluciones. La numeración de temas, secciones, enunciados y ecuaciones es la misma que en los apuntes, de modo que cualquier referencia puede buscarse allí; las que remiten a elementos suprimidos (un ejemplo, una figura, una ecuación de una demostración) conservan su número en los apuntes.

\section*{Notación}
\begin{center}
\begin{tabular}{@{}ll@{}}
\toprule
Símbolo & Significado \\
\midrule
$n$, $d$ & longitud de la secuencia y dimensión del modelo \\
$x_i\in\R^{1\times d}$, $X\in\R^{n\times d}$ & vector (fila) del token $i$ y matriz que los apila \\
$q_i,k_j\in\R^{1\times d_k}$, $v_j\in\R^{1\times d_v}$ & consulta (\emph{query}), clave (\emph{key}) y valor (\emph{value}) \\
$Q,K,V$ & matrices que apilan consultas, claves y valores \\
$\WQ,\WK\in\R^{d\times d_k}$, $\WV\in\R^{d\times d_v}$ & matrices de proyección aprendidas \\
$S=QK\T$, $s_{ij}=q_i\cdot k_j$ & matriz de puntuaciones (parecidos) \\
$A$, $\alpha_{ij}$ & matriz de atención y peso que la posición $i$ asigna a la $j$ \\
$O=AV$, $o_i$ & salida de la atención \\
$M$ & máscara aditiva, con entradas $0$ o $-\infty$ \\
$h$, $\WO$ & número de cabezas y proyección de salida \\
$\beta$, $\tau=1/\beta$ & inverso de la temperatura y temperatura \\
$P$ & matriz de permutación \\
$\conv\{v_1,\dots,v_n\}$ & envolvente convexa de los valores \\
$e_t$, $p_i$ & vectores \emph{one-hot} de token y de posición (tema~5) \\
\bottomrule
\end{tabular}
\end{center}
Los vectores de token, consulta, clave y valor son vectores fila, como en~\cite{vaswani} y~\cite{jm3}; así, $Q=X\WQ$. La softmax de una matriz se aplica por filas.


\tableofcontents

\mainmatter
% =====================================================================
\chapter{Atención como agrupación ponderada}
\label{tema:1}
% =====================================================================

Este tema define la atención en su forma más general: una consulta se compara con un conjunto de claves, los parecidos se convierten en pesos mediante una softmax y la salida es la media de los valores asociados ponderada por esos pesos. Se demuestran las dos propiedades estructurales básicas de la operación, la pertenencia de la salida a la envolvente convexa de los valores y su comportamiento en los límites de temperatura, y se presentan sus dos antecedentes: la regresión por núcleos de Nadaraya y Watson y la atención aditiva de Bahdanau, Cho y Bengio. Se siguen las §11.1, §11.2 y §11.4 de~\cite{d2l}, la §3.1 de~\cite{bahdanau}, la §7.1 de~\cite{jm3} y la clase~5 de~\cite{cs224n}.

\section{Del diccionario a la búsqueda suave}

Un diccionario almacena pares $(k_j,v_j)$ y, ante una consulta $q$, devuelve el valor $v_j$ cuya clave coincide con $q$, o un error si ninguna coincide. La búsqueda es \emph{dura}: el resultado depende de $q$ de forma discontinua y no admite gradientes. La atención sustituye la coincidencia exacta por un grado de parecido y la elección de un valor por una media ponderada~\cite[§11.1]{d2l}.

\begin{definicion}[Agrupación por atención]\label{def:agrupacion}
Sean $q\in\R^{1\times d_k}$ una \emph{consulta}, $k_1,\dots,k_n\in\R^{1\times d_k}$ unas \emph{claves} y $v_1,\dots,v_n\in\R^{1\times d_v}$ los \emph{valores} asociados. Sea $a\colon\R^{1\times d_k}\times\R^{1\times d_k}\to\R$ una \emph{función de puntuación}. La \emph{agrupación por atención} de los valores respecto de $q$ es
\[
\Att\bigl(q;(k_j,v_j)_{j=1}^{n}\bigr)=\sum_{j=1}^{n}\alpha_j\,v_j,\qquad
\alpha_j=\frac{\exp a(q,k_j)}{\sum_{l=1}^{n}\exp a(q,k_l)},
\]
es decir, $(\alpha_1,\dots,\alpha_n)=\softmax\bigl(a(q,k_1),\dots,a(q,k_n)\bigr)$. Los números $\alpha_j$ se llaman \emph{pesos de atención}.
\end{definicion}

La función de puntuación más usada es el producto escalar, $a(q,k)=q\cdot k=qk\T$, que es grande cuando los dos vectores apuntan en direcciones parecidas y tienen norma grande (módulo~2). Las claves deciden \emph{cuánto} contribuye cada par; los valores, \emph{qué} se contribuye. Separar ambos papeles permite que el criterio por el que un elemento es encontrado no coincida con la información que entrega.

\addtocounter{figure}{1}%

\begin{definicion}[Autoatención]\label{def:autoatencion}
Sean $x_1,\dots,x_n\in\R^{1\times d}$ los vectores de los tokens de una secuencia y $\WQ,\WK\in\R^{d\times d_k}$, $\WV\in\R^{d\times d_v}$ matrices de parámetros. La \emph{autoatención} asigna a cada posición $i$ la salida $o_i=\Att\bigl(q_i;(k_j,v_j)_{j=1}^{n}\bigr)$, con
\[
q_i=x_i\WQ,\qquad k_j=x_j\WK,\qquad v_j=x_j\WV .
\]
Si las consultas proceden de una secuencia y las claves y valores de otra, se habla de \emph{atención cruzada}.
\end{definicion}

\begin{observacion}
La \cref{fig:qkv} resume la definición. En un modelo de lenguaje del tipo \emph{decoder-only}, como los que se estudian en este curso, toda la atención es autoatención; la atención cruzada aparece en los modelos codificador-decodificador de traducción (\cref{sec:bahdanau}).
\end{observacion}

\section{Propiedades de la agrupación}

\begin{proposicion}[Envolvente convexa]\label{prop:convexa}
Para cualesquiera $q$, claves, valores y función de puntuación, la salida de la \cref{def:agrupacion} pertenece a la envolvente convexa $\conv\{v_1,\dots,v_n\}$. En particular, para cada coordenada $c$,
\[
\min_j v_j[c]\;\le\;\Att(q)[c]\;\le\;\max_j v_j[c].
\]
\end{proposicion}

La atención, por tanto, solo \emph{mezcla} información existente: no puede producir una coordenada mayor que la mayor de los valores. Por esa razón un bloque Transformer combina la atención con una red de alimentación hacia delante y con conexiones residuales (módulo~14).

Para estudiar la dependencia respecto de la escala de las puntuaciones se introduce un factor $\beta>0$, el \emph{inverso de la temperatura} del módulo~1: $\alpha_j(\beta)=\exp(\beta s_j)/\sum_l\exp(\beta s_l)$, con $s_j=a(q,k_j)$.

\begin{teorema}[Límites de temperatura]\label{teo:limites}
Sean $s_1,\dots,s_n$ las puntuaciones y supóngase que el máximo $s^{\ast}=\max_js_j$ se alcanza en un único índice $j^{\ast}$. Entonces
\[
\lim_{\beta\to+\infty}\Att_\beta(q)=v_{j^{\ast}},\qquad
\lim_{\beta\to 0^{+}}\Att_\beta(q)=\frac{1}{n}\sum_{j=1}^{n}v_j .
\]
\end{teorema}

La atención interpola, pues, entre el diccionario exacto (límite $\beta\to\infty$, que devuelve el valor de la clave más parecida) y la media simple de los valores ($\beta\to 0$). El caso con empates se trata en el \cref{eje:1-limites}.

\section{Regresión por núcleos de Nadaraya y Watson}

La idea de promediar respuestas con pesos que dependen del parecido entre la consulta y unos puntos de referencia es anterior al aprendizaje profundo. Zhang \emph{et al.}~\cite[§11.2]{d2l} la presentan a partir de la regresión no paramétrica.

\begin{definicion}[Estimador de Nadaraya--Watson]\label{def:nw}
Sean $(k_j,v_j)_{j=1}^{n}$ observaciones de entrada y salida y $\kappa(q,k)\ge 0$ un \emph{núcleo}. El estimador de Nadaraya--Watson en el punto $q$ es
\[
\hat f(q)=\sum_{j=1}^{n}\frac{\kappa(q,k_j)}{\sum_{l}\kappa(q,k_l)}\,v_j .
\]
Con el núcleo gaussiano, $\kappa(q,k)=\exp\bigl(-\tfrac12\norm{q-k}^{2}\bigr)$.
\end{definicion}

\begin{proposicion}\label{prop:gauss}
Con el núcleo gaussiano, el estimador de Nadaraya--Watson es la agrupación por atención con puntuación $a(q,k)=-\tfrac12\norm{q-k}^{2}$. Si además todas las claves tienen la misma norma, coincide con la agrupación con puntuación $a(q,k)=q\cdot k$.
\end{proposicion}

\section{La atención aditiva de Bahdanau}
\label{sec:bahdanau}

En un traductor codificador-decodificador con redes recurrentes, el codificador produce un vector $h_j$ por cada palabra de la frase de origen. En el modelo clásico, el decodificador solo recibía el último de ellos, un resumen de tamaño fijo de toda la frase; Bahdanau, Cho y Bengio propusieron que en cada paso el decodificador construyese su propio resumen como media ponderada de todos los $h_j$~\cite[§3.1]{bahdanau}.

\begin{definicion}[Atención aditiva]\label{def:aditiva}
Sean $h_1,\dots,h_{T}$ los vectores del codificador y $s_{i-1}$ el estado del decodificador antes de emitir la palabra $i$. El \emph{vector de contexto} es
\[
c_i=\sum_{j=1}^{T}\alpha_{ij}h_j,\qquad \alpha_{ij}=\frac{\exp e_{ij}}{\sum_{l}\exp e_{il}},\qquad e_{ij}=v_a\T\tanh\bigl(W_as_{i-1}+U_ah_j\bigr),
\]
donde $W_a$, $U_a$ y $v_a$ son parámetros aprendidos (los vectores de esta definición son columnas, como en el artículo original).
\end{definicion}

Es una agrupación por atención en la que la consulta es $s_{i-1}$, claves y valores son los $h_j$ y la puntuación es una red de una capa oculta~\cite[apéndice~A.1.2]{bahdanau}. Vaswani \emph{et al.} señalan que la puntuación por producto escalar tiene una complejidad teórica similar, pero es mucho más rápida en la práctica porque se reduce a productos de matrices optimizados~\cite[§3.2.1]{vaswani}; es la que se estudia en el resto del módulo.


% =====================================================================
\chapter{Atención por producto escalar escalado}
\label{tema:2}
% =====================================================================

Este tema escribe la autoatención de todas las posiciones a la vez como dos productos de matrices y una softmax por filas, resuelve a mano un ejemplo completo, demuestra que la operación es equivariante por permutaciones y justifica el factor $1/\sqrt{d_k}$ con un argumento de varianza y con la jacobiana de la softmax. Se siguen la §3.2.1 de~\cite{vaswani}, las §7.1.1 y §7.3 de~\cite{jm3}, la §11.3 de~\cite{d2l} y la §12.1 de~\cite{bishopdl}.

\section{Forma matricial}

Se apilan las consultas, claves y valores de las $n$ posiciones como filas de $Q=X\WQ\in\R^{n\times d_k}$, $K=X\WK\in\R^{n\times d_k}$ y $V=X\WV\in\R^{n\times d_v}$.

\begin{definicion}[Atención por producto escalar escalado]\label{def:sdpa}
Para $Q\in\R^{n\times d_k}$, $K\in\R^{m\times d_k}$ y $V\in\R^{m\times d_v}$,
\[
\Att(Q,K,V)=\softmax\!\Bigl(\frac{QK\T}{\sqrt{d_k}}\Bigr)V,
\]
donde la softmax se aplica a cada fila. La matriz $A=\softmax(QK\T/\sqrt{d_k})\in\R^{n\times m}$ se llama \emph{matriz} (o \emph{mapa}) \emph{de atención}.
\end{definicion}

La entrada $(i,j)$ de $S=QK\T$ es $q_i\cdot k_j$, de modo que la fila $i$ de $A$ contiene los pesos de la \cref{def:agrupacion} para la consulta $q_i$ con puntuación $q\cdot k/\sqrt{d_k}$, y la fila $i$ de $AV$ es la salida $o_i$. La \cref{tab:formas} resume las formas. Vaswani \emph{et al.} introdujeron el nombre y la fórmula~\cite[§3.2.1, ec.~(1)]{vaswani}; la tarea~1 de CS336 la adopta con $m\ne n$ para cubrir la atención cruzada~\cite[§3.4.4]{cs336}.

\begin{table}[h]
\centering
\begin{tabular}{@{}clll@{}}
\toprule
Paso & Operación & Forma & Significado \\
\midrule
1 & $S=QK\T$ & $n\times m$ & $s_{ij}=q_i\cdot k_j$ \\
2 & $S/\sqrt{d_k}$ & $n\times m$ & escalado (\cref{sec:escalado}) \\
3 & $A=\softmax(S/\sqrt{d_k})$ & $n\times m$ & cada fila es una distribución \\
4 & $O=AV$ & $n\times d_v$ & fila $i$: media ponderada de valores \\
\bottomrule
\end{tabular}
\caption{Formas de las matrices intermedias de la atención.}
\label{tab:formas}
\end{table}

\addtocounter{definicion}{1}%

\addtocounter{figure}{1}%

\section{Equivariancia por permutaciones}

\begin{proposicion}[Equivariancia]\label{prop:equivariancia}
Sea $f(X)=\softmax\bigl(X\WQ(X\WK)\T/\sqrt{d_k}\bigr)X\WV$ la autoatención sin máscara y sea $P\in\R^{n\times n}$ una matriz de permutación. Entonces $f(PX)=P\,f(X)$.
\end{proposicion}

La autoatención no distingue, por sí sola, el orden de los tokens: «el perro muerde al hombre» y «el hombre muerde al perro» producen las mismas salidas, solo que en otras filas. Por eso los modelos de lenguaje añaden información de posición a los vectores de entrada (módulo~14), y por eso la máscara causal del \cref{tema:3}, que sí depende del orden, rompe esta simetría.

\section{Por qué se divide por \texorpdfstring{$\sqrt{d_k}$}{raíz de dk}}
\label{sec:escalado}

Vaswani \emph{et al.} justifican el factor con un argumento de varianza que enuncian en una nota al pie~\cite[§3.2.1, nota~4]{vaswani}; Zhang \emph{et al.}~\cite[§11.3.3]{d2l} y Jurafsky y Martin~\cite[§7.1.1]{jm3} lo reproducen. Se da aquí con demostración.

\begin{teorema}[Varianza del producto escalar]\label{teo:varianza}
Sean $q=(q_1,\dots,q_{d_k})$ y $k=(k_1,\dots,k_{d_k})$ vectores aleatorios cuyas $2d_k$ componentes son independientes, de media~0 y varianza~1. Entonces
\[
\E[q\cdot k]=0,\qquad \Var(q\cdot k)=d_k,\qquad \Var\Bigl(\frac{q\cdot k}{\sqrt{d_k}}\Bigr)=1 .
\]
\end{teorema}

Por el teorema central del límite, $q\cdot k/\sqrt{d_k}$ es aproximadamente $\mathcal N(0,1)$ para $d_k$ grande. Sin el factor, con $d_k=64$ los parecidos tendrían desviación típica~8 en la inicialización y diferencias de~16 entre dos claves serían frecuentes. Una simulación con $2\cdot 10^{5}$ pares de vectores gaussianos da varianzas empíricas $4{,}00$ para $d_k=4$ y $63{,}8$ para $d_k=64$.

Para ver por qué importa la escala se necesita la jacobiana de la softmax, calculada en el módulo~4 con la regla del cociente.

\begin{proposicion}[Saturación de la softmax]\label{prop:saturacion}
Sea $p=\softmax(z)\in\R^{n}$ y $J=\partial p/\partial z$, de entradas $J_{ij}=p_i(\delta_{ij}-p_j)$. Entonces
\[
\max_{i,j}\abs{J_{ij}}\;\le\;1-\max_{l}p_l .
\]
En consecuencia, si una probabilidad tiende a~1, todas las derivadas parciales de la softmax tienden a~0.
\end{proposicion}

\begin{table}[h]
\centering
\begin{tabular}{@{}lll@{}}
\toprule
Logits $z$ & $\softmax(z)$ & $\max\abs{J_{ij}}$ \\
\midrule
$(1,0,0)$ & $(0{,}576;\allowbreak\ 0{,}212;\allowbreak\ 0{,}212)$ & $0{,}244$ \\
$(8,0,0)$ & $(0{,}9993;\allowbreak\ 0{,}0003;\allowbreak\ 0{,}0003)$ & $0{,}00067$ \\
\bottomrule
\end{tabular}
\caption{Con logits ocho veces mayores, las derivadas de la softmax se reducen en un factor de unos 360.}
\label{tab:saturacion}
\end{table}

\begin{observacion}[El factor como temperatura]
Dividir por $\sqrt{d_k}$ equivale a aplicar la softmax con temperatura $\tau=\sqrt{d_k}$ (módulo~1, tema~4). Sin ella, el modelo empezaría en la región saturada: en el \cref{ej:tres}, sin dividir por~2, el peso de la fila~2 sobre la posición~2 pasa de $0{,}867$ a $0{,}982$. Vaswani \emph{et al.} observan, en la misma línea, que sin escalado la atención aditiva supera a la de producto escalar cuando $d_k$ es grande~\cite[§3.2.1]{vaswani}. Durante el entrenamiento, la red puede modificar la escala efectiva agrandando o encogiendo $\WQ$ y $\WK$.
\end{observacion}

\begin{observacion}[Evaluación estable]
Como la softmax es invariante al restar una constante a todas las entradas, se evalúa restando el máximo de cada fila antes de exponenciar, lo que evita desbordamientos en fp16 (módulo~11). La tarea~1 de CS336 exige implementarla así~\cite[§3.4.4]{cs336}.
\end{observacion}


% =====================================================================
\chapter{La máscara causal}
\label{tema:3}
% =====================================================================

Un modelo de lenguaje autorregresivo factoriza la probabilidad de una secuencia con la regla de la cadena (módulo~1, tema~2): la distribución del token $t+1$ solo puede depender de los tokens $1,\dots,t$. Este tema define las máscaras aditivas, demuestra que sumar $-\infty$ antes de la softmax equivale a restringir la softmax al conjunto de posiciones permitidas, prueba que la máscara causal hace que cada salida dependa solo de su pasado y deduce de ello que una única pasada produce las $n$ predicciones de una secuencia. Se siguen la §3.2.3 de~\cite{vaswani}, la §7.3 de~\cite{jm3}, las §3.4.4 y §3.4.5 de la tarea~1 de CS336~\cite{cs336} y la §11.3.2.1 de~\cite{d2l}.

\section{Máscaras aditivas}

\begin{definicion}[Máscara aditiva]\label{def:mascara}
Sea $\mathcal{P}\subseteq\{1,\dots,n\}\times\{1,\dots,m\}$ el conjunto de pares $(i,j)$ \emph{permitidos}, con la condición de que cada fila $i$ contenga al menos un par permitido. La \emph{máscara aditiva} asociada es la matriz $M\in\{0,-\infty\}^{n\times m}$ con $M_{ij}=0$ si $(i,j)\in\mathcal{P}$ y $M_{ij}=-\infty$ en otro caso. La \emph{atención enmascarada} es
\[
\Att_M(Q,K,V)=\softmax\!\Bigl(\frac{QK\T}{\sqrt{d_k}}+M\Bigr)V,
\]
con el convenio $e^{-\infty}=0$.
\end{definicion}

\begin{definicion}[Máscara causal]\label{def:causal}
En la autoatención ($m=n$), la \emph{máscara causal} es la asociada a $\mathcal{P}=\{(i,j):j\le i\}$:
\[
M_{ij}=\begin{cases}0,& j\le i,\\-\infty,& j>i.\end{cases}
\]
Es triangular estrictamente superior en sus entradas $-\infty$; la diagonal siempre está permitida.
\end{definicion}

Vaswani \emph{et al.} la describen como el bloqueo de las conexiones «ilegales» de la autoatención del decodificador, fijando a $-\infty$ las entradas correspondientes de la softmax~\cite[§3.2.3]{vaswani}; Jurafsky y Martin la presentan en la §7.3 de~\cite{jm3}, y la tarea~1 de CS336 la implementa con una máscara booleana en la que \texttt{True} indica «puede atender»~\cite[§3.4.4]{cs336}. La \cref{fig:causal} muestra su estructura.

\addtocounter{figure}{1}%

\begin{proposicion}[Softmax restringida]\label{prop:restringida}
Sea $z\in\R^{m}$ y sea $\mathcal{P}_i\subseteq\{1,\dots,m\}$ no vacío. Si $\mu_j=0$ para $j\in\mathcal{P}_i$ y $\mu_j=-\infty$ en otro caso, entonces
\[
\softmax(z+\mu)_j=\begin{cases}\dfrac{e^{z_j}}{\sum_{l\in\mathcal{P}_i}e^{z_l}},& j\in\mathcal{P}_i,\\[2ex] 0,& j\notin\mathcal{P}_i.\end{cases}
\]
En particular, los pesos permitidos suman~1 y no dependen de las puntuaciones $z_j$ con $j\notin\mathcal{P}_i$.
\end{proposicion}

La condición $\mathcal{P}_i\ne\varnothing$ es necesaria: una fila enteramente enmascarada daría $0/0$, que en coma flotante produce \texttt{NaN} (módulo~11). En la máscara causal no puede ocurrir, porque $(i,i)$ siempre está permitido.

\begin{observacion}[Por qué $-\infty$ antes de la softmax]
Sumar~0 antes de la softmax no anula nada, porque $e^{0}=1$. Poner a~0 los pesos \emph{después} de la softmax deja filas que suman menos de~1 y pesos del pasado más pequeños de lo debido; hay que renormalizar, y el resultado coincide entonces con el de la \cref{prop:restringida} (\cref{eje:3-renorm}). En la práctica se usa \texttt{-inf} o un número negativo grande; el \cref{eje:3-finito} cuantifica cuánto debe valer.
\end{observacion}

\addtocounter{definicion}{1}%

\section{Causalidad}

\begin{teorema}[Causalidad de la atención enmascarada]\label{teo:causal}
Sean $q_i=x_i\WQ$, $k_j=x_j\WK$, $v_j=x_j\WV$ y sea $o_i$ la fila $i$ de la autoatención con máscara causal. Entonces, para cada $i$, $o_i$ es función únicamente de $x_1,\dots,x_i$.
\end{teorema}

\begin{corolario}[Entrenamiento en paralelo]\label{cor:paralelo}
Sea $f$ la autoatención con máscara causal y, para cada $t\le n$, sea $f_t$ la misma operación aplicada solo al prefijo $x_{1:t}$. Entonces la fila $t$ de $f(x_{1:n})$ coincide con la última fila de $f_t(x_{1:t})$. En consecuencia, una sola evaluación sobre la secuencia completa proporciona las $n$ representaciones que una red usa para predecir $x_2,\dots,x_{n+1}$, cada una calculada solo con su prefijo.
\end{corolario}

La tarea~1 de CS336 describe la alternativa ingenua, ejecutar la atención $n$ veces, una por prefijo, y la descarta en favor de la máscara~\cite[§3.4.5]{cs336}. Una red recurrente necesita, en cambio, $n$ pasos secuenciales, porque el estado en $t$ depende del estado en $t-1$ (módulo~12). Esta diferencia, que Vaswani \emph{et al.} cuantifican como $O(1)$ operaciones secuenciales frente a $O(n)$~\cite[§4, tabla~1]{vaswani}, es la razón práctica principal por la que los Transformer desplazaron a las redes recurrentes en el preentrenamiento de modelos de lenguaje.

\begin{observacion}[Otras máscaras]
La \cref{def:mascara} admite otros conjuntos $\mathcal{P}$. La máscara de relleno (\emph{padding}) prohíbe las posiciones añadidas para igualar la longitud de las secuencias de un lote~\cite[§11.3.2.1]{d2l}; la de ventana deslizante permite solo $i-w<j\le i$, lo que reduce el coste a $O(nw)$ (\cref{tema:4}). La combinación de varias máscaras es la intersección de sus conjuntos permitidos, es decir, la suma de sus matrices.
\end{observacion}


% =====================================================================
\chapter{Atención multi-cabeza y coste computacional}
\label{tema:4}
% =====================================================================

Una sola cabeza de atención define una única noción de parecido entre posiciones. La atención multi-cabeza ejecuta varias en paralelo, cada una en un subespacio de dimensión $d/h$, y combina sus resultados con una proyección lineal. Este tema la define, demuestra que su número de parámetros no depende del número de cabezas, muestra que cada cabeza tiene una forma bilineal de rango acotado, cuenta las operaciones de coma flotante de la capa y estudia la memoria de la matriz de atención, que es el obstáculo principal para los contextos largos. Se siguen la §3.2.2 y la §4 de~\cite{vaswani}, las §3.4.5 y §3.5 de la tarea~1 de CS336~\cite{cs336}, las §11.5 y §11.6.2 de~\cite{d2l} y las §7.1.1 y §7.3 de~\cite{jm3}.

\section{Definición}

\begin{definicion}[Atención multi-cabeza]\label{def:mha}
Sean $d$ la dimensión del modelo, $h$ un divisor de $d$ y $d_k=d_v=d/h$. Para cada cabeza $r=1,\dots,h$ se toman matrices $\WQ^{(r)},\WK^{(r)},\WV^{(r)}\in\R^{d\times d/h}$, y además $\WO\in\R^{d\times d}$. La \emph{atención multi-cabeza} de $X\in\R^{n\times d}$ es
\[
\MHA(X)=\Concat\bigl(H_1,\dots,H_h\bigr)\,\WO,\qquad
H_r=\Att\bigl(X\WQ^{(r)},X\WK^{(r)},X\WV^{(r)}\bigr)\in\R^{n\times d/h},
\]
donde $\Concat$ yuxtapone las $h$ matrices por columnas y la atención de cada cabeza puede llevar máscara causal.
\end{definicion}

Vaswani \emph{et al.} usan $d=512$ y $h=8$, de modo que $d_k=d_v=64$, y observan que el coste total es similar al de una sola cabeza de dimensión completa~\cite[§3.2.2]{vaswani}. Zhang \emph{et al.} dan la definición con dimensiones generales y la misma elección~\cite[§11.5.1]{d2l}.

\addtocounter{figure}{1}%

\begin{observacion}[Implementación]
No se calculan $h$ productos pequeños. Se apilan las matrices de las cabezas en $\WQ=\bigl(\WQ^{(1)}\ \cdots\ \WQ^{(h)}\bigr)\in\R^{d\times d}$, se calcula $Q=X\WQ$ con un único producto y se reordena el resultado de forma $(n,d)$ a $(h,n,d/h)$; la fórmula de la atención se aplica entonces con un eje adicional de cabezas, que se trata como un eje de lote (módulo~3). La tarea~1 de CS336 pide calcular así las tres proyecciones con tres productos de matrices, o con uno solo concatenando las tres matrices~\cite[§3.4.5]{cs336}.
\end{observacion}

\section{Parámetros y estructura}

\begin{proposicion}[Parámetros]\label{prop:param}
Sin sesgos, la atención multi-cabeza tiene $4d^{2}$ parámetros, cualquiera que sea $h$.
\end{proposicion}

Con $d=768$, la capa tiene $4\cdot 768^{2}=2\,359\,296$ parámetros; con $d=512$, $1\,048\,576$. En el módulo~14 se suman los $8d^{2}$ de la red de alimentación hacia delante con $d_{ff}=4d$ y se obtiene que la atención es un tercio de los parámetros de cada bloque.

\begin{proposicion}[Descomposición por cabezas]\label{prop:descomp}
Sea $\WO^{(r)}\in\R^{d/h\times d}$ el bloque de filas $(r-1)d/h+1,\dots,rd/h$ de $\WO$. Entonces
\[
\MHA(X)=\sum_{r=1}^{h}H_r\,\WO^{(r)}=\sum_{r=1}^{h}A_r\,X\,\WV^{(r)}\WO^{(r)},
\]
donde $A_r$ es la matriz de atención de la cabeza $r$. Además, las puntuaciones de la cabeza $r$ son $s^{(r)}_{ij}=x_i\,\WQ^{(r)}\WK^{(r)\top}x_j\T$, una forma bilineal cuya matriz $W_{QK}^{(r)}=\WQ^{(r)}\WK^{(r)\top}\in\R^{d\times d}$ tiene rango a lo sumo $d/h$.
\end{proposicion}

La proposición muestra que la multi-cabeza no es una sola atención grande cortada en trozos: es una \emph{suma} de $h$ atenciones, cada una con su propia matriz de pesos $A_r$, una forma bilineal de rango $\le d/h$ que decide a dónde mirar y una aplicación $\WV^{(r)}\WO^{(r)}$ que decide qué escribir. Esta lectura es la base del análisis de circuitos del \cref{tema:5}. La tarea~3 de CS224N ilustra con un ejemplo de claves ruidosas por qué promediar dos cabezas puede ser más robusto que una sola~\cite[tarea~3, pregunta~1]{cs224n}.

\begin{observacion}[Variantes]
Muchos modelos abiertos comparten las claves y los valores entre grupos de cabezas (atención \emph{multi-query} y \emph{grouped-query}) para reducir la memoria de inferencia; se estudian en el módulo~18. Qué variante usan Claude, Gemini o Grok en sus versiones actuales no se ha publicado en detalle.
\end{observacion}

\section{Coste en operaciones}

Se cuenta en operaciones de coma flotante (FLOP), con el convenio de CS336 de que una multiplicación seguida de una suma cuenta como dos~\cite[§3.5]{cs336}.

\begin{lema}\label{lem:flops}
El producto de una matriz $a\times b$ por una $b\times c$ requiere $2abc$ FLOP.
\end{lema}

\begin{proposicion}[Coste de la capa]\label{prop:coste}
Con $h$ cabezas de dimensión $d/h$, una pasada de la atención multi-cabeza sobre $n$ posiciones requiere
\begin{align*}
\FLOP_{\text{proy}}&=8nd^{2} &&\text{(proyecciones $Q,K,V$ y $\WO$)},\\
\FLOP_{\text{at}}&=4n^{2}d &&\text{(productos $QK\T$ y $AV$)},
\end{align*}
sin contar la softmax, que es de orden $hn^{2}$. La parte cuadrática domina si y solo si $n>2d$.
\end{proposicion}

Con $d=512$, la atención propiamente dicha domina a partir de $1024$ tokens. Vaswani \emph{et al.} resumen esta comparación en su tabla~1: complejidad por capa $O(n^{2}d)$ para la autoatención frente a $O(nd^{2})$ para una red recurrente, a cambio de $O(1)$ operaciones secuenciales y longitud máxima de camino $O(1)$ entre dos posiciones~\cite[§4]{vaswani}. Zhang \emph{et al.} hacen el mismo análisis en~\cite[§11.6.2]{d2l}.

\section{Memoria y contexto largo}

La matriz de atención tiene $n^{2}$ entradas por cabeza y por capa. En fp16, con dos bytes por número (módulo~11), ocupa $2n^{2}$ bytes.

\begin{table}[h]
\centering
\begin{tabular}{@{}rr@{}}
\toprule
$n$ (tokens) & Memoria de $A$ (una cabeza, una capa, fp16) \\
\midrule
$4096$ & $0{,}034$ GB \\
$10^{5}$ & $20$ GB \\
$10^{6}$ & $2000$ GB \\
\bottomrule
\end{tabular}
\caption{Memoria de una matriz de atención completa ($1\text{ GB}=10^{9}$ bytes).}
\label{tab:memoria}
\end{table}

Multiplicar el contexto por $k$ multiplica esta memoria y los $4n^{2}d$ FLOP por $k^{2}$. Las técnicas que lo mitigan se estudian en módulos posteriores: FlashAttention~\cite{dao} calcula la softmax por bloques sin almacenar $A$ completa, de modo que la memoria adicional pasa a ser lineal en $n$ mientras el número de FLOP sigue siendo $O(n^{2}d)$ (módulo~17); las ventanas deslizantes y la atención dispersa reducen el conjunto permitido $\mathcal{P}$; y en inferencia la caché de claves y valores evita recalcular las del pasado (módulo~19).

\begin{observacion}[Contexto largo en modelos comerciales]
El informe técnico de Gemini~1.5~\cite{gemini15} documenta recuperación casi perfecta en experimentos con contextos de hasta diez millones de tokens, y Anthropic ofrece para Claude ventanas de cientos de miles de tokens. Ninguna de las dos empresas ha publicado cómo resuelven internamente el coste cuadrático.
\end{observacion}


% =====================================================================
\chapter{Circuitos de atención y memorias asociativas}
\label{tema:5}
% =====================================================================

Los temas anteriores tratan la atención como una operación aislada. Este tema estudia dos lecturas estructurales de ella. La primera, debida a Elhage \emph{et al.}~\cite{elhage}, descompone cada cabeza en un circuito que decide a dónde mirar y otro que decide qué copiar, y muestra que dos cabezas en capas sucesivas pueden componerse para implementar una \emph{cabeza de inducción}, el mecanismo que Olsson \emph{et al.}~\cite{olsson} relacionaron con el aprendizaje en contexto. La segunda, debida a Ramsauer \emph{et al.}~\cite{ramsauer}, identifica la atención con la regla de actualización de una red de Hopfield moderna, una memoria asociativa direccionable por contenido. Se siguen esos tres artículos y la §10.6 de~\cite{jm3}.

\section{El flujo residual y los circuitos QK y OV}

En un Transformer, cada capa \emph{suma} su salida a la representación de cada posición, que recorre así la red como un \emph{flujo residual} (módulos~11 y~14). Elhage \emph{et al.} estudian modelos formados solo por capas de atención, en los que el vector final de la posición $i$ es su embedding más las contribuciones de todas las cabezas de todas las capas~\cite{elhage}. Por la \cref{prop:descomp}, la contribución de una cabeza a la posición $i$ es
\[
\sum_{j\le i}\alpha_{ij}\;x_j\,W_{OV},\qquad
\alpha_{ij}=\softmax_j\Bigl(\frac{x_i\,W_{QK}\,x_j\T}{\sqrt{d_k}}\Bigr),
\]
con $W_{QK}=\WQ\WK\T$ y $W_{OV}=\WV\WO^{(r)}$, ambas de tamaño $d\times d$ (en esta sección se omite el índice de la cabeza).

\begin{definicion}[Circuitos QK y OV]\label{def:circuitos}
En una cabeza de atención, el \emph{circuito QK} es la forma bilineal $W_{QK}=\WQ\WK\T$, que determina los pesos de atención, y el \emph{circuito OV} es la aplicación lineal $W_{OV}=\WV\WO$, que determina qué se escribe en el flujo residual de la posición que atiende a partir de la posición atendida.
\end{definicion}

Los dos circuitos son independientes: el primero decide \emph{a dónde} se mira; el segundo, \emph{qué} se copia. Si los vectores de entrada son \emph{one-hot}, $x_iW_{QK}x_j\T$ es una entrada de la matriz $W_{QK}$ y la tabla completa de preferencias de la cabeza puede leerse directamente.

\section{Cabezas de inducción}

Una \emph{cabeza de inducción} implementa la regla «si antes apareció $a$ seguido de $b$ y ahora vuelve a aparecer $a$, predecir $b$», es decir, completa patrones del tipo $[a][b]\dots[a]\to[b]$~\cite{olsson}. Permite copiar del propio contexto secuencias que nunca se vieron en el entrenamiento, como un nombre propio inventado que se repite en un texto.

\begin{proposicion}\label{prop:una-capa}
En un modelo de una sola capa de atención cuyas entradas en cada posición son el embedding de su token más el de su posición, la puntuación $s_{ij}$ es función de $(x_i,x_j)$. En consecuencia, para dos posiciones permitidas $j,l\le i$, el cociente $\alpha_{ij}/\alpha_{il}$ solo depende de $x_i$, $x_j$ y $x_l$; en particular, ninguna cabeza de esa capa puede preferir la posición $j$ a la $l$ en función del token \emph{anterior} $x_{j-1}$.
\end{proposicion}

Hace falta, por tanto, que una capa anterior haya escrito en la posición $j$ información sobre $x_{j-1}$. La construcción de Elhage \emph{et al.} usa dos cabezas~\cite{elhage}:
\begin{enumerate}[label=\arabic*.]
  \item una \emph{cabeza de token anterior} en la capa~1, cuyo mapa de atención es la subdiagonal: la posición $j$ atiende a $j-1$ y su circuito OV copia la identidad del token, de modo que escribe en $j$ «mi token anterior era $x_{j-1}$»;
  \item una \emph{cabeza de inducción} en la capa~2, cuyas claves se construyen a partir de lo escrito por la capa~1 y cuya consulta es el token actual: la posición $i$ atiende a las $j$ con $x_{j-1}=x_i$, y su circuito OV copia el token $x_j$.
\end{enumerate}
Que las claves de la capa~2 lean la salida de una cabeza de la capa~1 se llama \emph{composición K}; análogamente se definen las composiciones Q y V. La \cref{fig:induccion} ilustra el mecanismo.

\addtocounter{figure}{1}%

\begin{teorema}[Cabeza de token anterior]\label{teo:previo}
Sean $p_0,\dots,p_{n-1}$ vectores \emph{one-hot} de posición y $\beta>0$. Con máscara causal, consultas $q_i=\beta p_i$ y claves $k_j=\beta p_{j+1}$ (con $p_n=0$), y sin escalado adicional, el peso que la posición $i\ge 1$ asigna a $i-1$ es
\[
\alpha_{i,i-1}=\frac{e^{\beta^{2}}}{e^{\beta^{2}}+i}\xrightarrow[\beta\to\infty]{}1 .
\]
\end{teorema}

Con valores $v_j=e_{x_j}$ (el \emph{one-hot} del token), la salida de esta cabeza en la posición $j$ es, para $\beta$ grande, aproximadamente $e_{x_{j-1}}$. El \cref{eje:5-induccion} completa la construcción con la segunda capa.

\begin{observacion}[Evidencia empírica]
Olsson \emph{et al.} presentan seis líneas de evidencia complementarias de que las cabezas de inducción explican una parte importante del aprendizaje en contexto: entre otras, que se forman durante el entrenamiento en el mismo intervalo en el que la capacidad de aprendizaje en contexto mejora bruscamente, visible como un escalón en la curva de pérdida. La evidencia es causal en modelos pequeños formados solo por atención y correlacional en modelos mayores con redes de alimentación hacia delante~\cite{olsson}. Jurafsky y Martin resumen estos resultados en la §10.6 de~\cite{jm3}; el módulo~23 retoma los circuitos.
\end{observacion}

\section{Atención y redes de Hopfield modernas}

Una memoria asociativa almacena patrones y, ante un estímulo parecido a uno de ellos, devuelve el patrón completo. Ramsauer \emph{et al.} definen una red de Hopfield con estados continuos cuya regla de actualización es la atención~\cite[§2]{ramsauer}. En esta sección se usa su convención: $X=(x_1,\dots,x_N)\in\R^{d\times N}$ tiene los patrones almacenados como \emph{columnas} y el estado $\xi\in\R^{d}$ es un vector columna.

\begin{definicion}[Energía de Hopfield moderna]\label{def:hopfield}
Sean $\beta>0$ y $\lse(\beta,z)=\beta^{-1}\ln\sum_{l=1}^{N}e^{\beta z_l}$. La energía del estado $\xi$ es
\[
E(\xi)=-\lse\bigl(\beta,X\T\xi\bigr)+\tfrac12\,\xi\T\xi+\beta^{-1}\ln N+\tfrac12M^{2},\qquad M=\max_l\norm{x_l},
\]
y la \emph{regla de actualización} es $\xi^{\text{nuevo}}=X\softmax\bigl(\beta X\T\xi\bigr)$.
\end{definicion}

\begin{proposicion}[Puntos estacionarios]\label{prop:hopfield}
El gradiente de la energía es $\nabla E(\xi)=\xi-X\softmax(\beta X\T\xi)$. Por tanto, $\xi$ es un punto estacionario de $E$ si y solo si es un punto fijo de la regla de actualización, y la actualización consiste en saltar al punto que anula el gradiente de la parte cuadrática con la parte $\lse$ congelada.
\end{proposicion}

Ramsauer \emph{et al.} demuestran además que la actualización hace decrecer la energía y converge, que la red puede almacenar un número de patrones exponencial en la dimensión y que, si los patrones están bien separados, un solo paso basta para recuperar uno de ellos con error exponencialmente pequeño~\cite[§2, teoremas~1 a~5]{ramsauer}.

\begin{observacion}[Hopfield como atención]
Si los patrones almacenados y los estados se proyectan linealmente para formar claves, valores y consultas y se toma $\beta=1/\sqrt{d_k}$, la regla de actualización, escrita con vectores fila, adopta la forma $\softmax(QK\T/\sqrt{d_k})V$ de la \cref{def:sdpa}~\cite[§2]{ramsauer}. La atención es, pues, un paso de recuperación en una memoria asociativa, y $\beta$ desempeña el papel de inverso de una temperatura: el \cref{teo:limites} describe sus dos regímenes extremos, recuperación de un patrón y media de todos ellos.
\end{observacion}



\backmatter
% =====================================================================
\begin{thebibliography}{99}
\addcontentsline{toc}{chapter}{Bibliografía}
\markboth{Bibliografía}{}
\small

\bibitem{bahdanau}
D.~Bahdanau, K.~Cho e Y.~Bengio.
\newblock Neural machine translation by jointly learning to align and translate.
\newblock En \emph{International Conference on Learning Representations (ICLR)}, 2015.
\newblock \url{https://arxiv.org/abs/1409.0473}

\bibitem{bishopdl}
C.~M. Bishop y H.~Bishop.
\newblock \emph{Deep Learning: Foundations and Concepts}.
\newblock Springer, 2024.
\newblock Cap.~12, «Transformers» (§12.1 «Attention»).
\newblock \url{https://www.bishopbook.com/}

\bibitem{cs224n}
D.~Yang y Y.~Choi.
\newblock CS224N: Natural Language Processing with Deep Learning.
\newblock Stanford University, invierno de 2026. Clase~5, «Transformers»; tarea~3, «Self-Attention and Transformers».
\newblock \url{https://web.stanford.edu/class/cs224n/}

\bibitem{cs336}
Equipo docente de CS336.
\newblock CS336: Language Modeling from Scratch. Tarea~1, «Basics», versión 26.0.3.
\newblock Stanford University, primavera de 2026.
\newblock \url{https://github.com/stanford-cs336/assignment1-basics}

\bibitem{dao}
T.~Dao, D.~Y. Fu, S.~Ermon, A.~Rudra y C.~Ré.
\newblock FlashAttention: Fast and memory-efficient exact attention with IO-awareness.
\newblock En \emph{Advances in Neural Information Processing Systems 35}, 2022.
\newblock \url{https://arxiv.org/abs/2205.14135}

\bibitem{elhage}
N.~Elhage, N.~Nanda, C.~Olsson \emph{et al.}
\newblock A mathematical framework for Transformer circuits.
\newblock \emph{Transformer Circuits Thread}, Anthropic, 2021.
\newblock \url{https://transformer-circuits.pub/2021/framework/index.html}

\bibitem{gemini15}
Gemini Team, Google.
\newblock Gemini 1.5: Unlocking multimodal understanding across millions of tokens of context.
\newblock Informe técnico, 2024.
\newblock \url{https://arxiv.org/abs/2403.05530}

\bibitem{jm3}
D.~Jurafsky y J.~H. Martin.
\newblock \emph{Speech and Language Processing}, 3.ª ed., borrador del 19 de agosto de 2026.
\newblock Cap.~7, «Transformers and Pretraining»; cap.~10, «Interpretability».
\newblock \url{https://web.stanford.edu/~jurafsky/slp3/}

\bibitem{olsson}
C.~Olsson, N.~Elhage, N.~Nanda \emph{et al.}
\newblock In-context learning and induction heads.
\newblock \emph{Transformer Circuits Thread}, Anthropic, 2022.
\newblock \url{https://arxiv.org/abs/2209.11895}

\bibitem{ramsauer}
H.~Ramsauer, B.~Schäfl, J.~Lehner, P.~Seidl, M.~Widrich, T.~Adler, L.~Gruber, M.~Holzleitner, M.~Pavlović, G.~K. Sandve, V.~Greiff, D.~Kreil, M.~Kopp, G.~Klambauer, J.~Brandstetter y S.~Hochreiter.
\newblock Hopfield networks is all you need.
\newblock En \emph{International Conference on Learning Representations (ICLR)}, 2021.
\newblock \url{https://arxiv.org/abs/2008.02217}

\bibitem{vaswani}
A.~Vaswani, N.~Shazeer, N.~Parmar, J.~Uszkoreit, L.~Jones, A.~N. Gomez, Ł.~Kaiser e I.~Polosukhin.
\newblock Attention is all you need.
\newblock En \emph{Advances in Neural Information Processing Systems 30}, 2017.
\newblock \url{https://arxiv.org/abs/1706.03762}

\bibitem{d2l}
A.~Zhang, Z.~C. Lipton, M.~Li y A.~J. Smola.
\newblock \emph{Dive into Deep Learning}, versión 1.0.3.
\newblock Cambridge University Press, 2023.
\newblock Cap.~11, «Attention Mechanisms and Transformers».
\newblock \url{https://d2l.ai/}

\end{thebibliography}


\end{document}
